\chapter{离线生成动力学：梦境、记忆回放与模型校准}
\label{chap:physics_of_dreams}
梦境涉及外部输入降低、内部生成增强、记忆重放、情绪加工和元认知改变等多种过程，不能由单一“升温”或“流形自由演化”概括。我们在本章把梦境作为一种受生理状态约束的离线生成特例：系统仍接收内感受和残余外部信号，行动通道受到门控，内部模型以不同于清醒任务的方式采样与组合。HSF--HD负责描述状态、输入、结构、控制和报告如何改变；神经机制与临床现象须服从相应领域证据，AgentOS的离线维护只能作为工程对照，不能反向证明人脑按同一模块运行。每项功能解释还要与无梦睡眠、清醒想象和离线训练等替代条件比较。

\section{输入与行动边界的重配置：从现实闭环到内源生成}
\label{sec:section_3039}
清醒与睡眠的差别不是从狄利克雷边界突然切换到诺伊曼边界，更不是从开放系统变成绝热封闭系统。睡眠中视觉、听觉和运动输出的增益会变化，但内感受、姿态、呼吸、温度、噪声和触觉仍可进入；梦境也会受近期事件与身体状态影响。我们把边界写成多通道增益与权限的联合状态，而不是把“现实”设为场在边界上的固定值。

令外部观测通道为 $u_k^{ext}\in\mathbb R^{d_e}$，内感受通道为 $u_k^{int}\in\mathbb R^{d_i}$，内部生成候选为 $z_k\in\mathbb R^{d_z}$，动作命令为 $a_k$。睡眠模式 $m_k$ 决定感知增益 $G_e(m_k)$、内感受增益 $G_i(m_k)$、动作门控 $G_a(m_k)$ 与生成增益 $G_z(m_k)$。工作态离散更新可写为
\[
h_{k+1}=F\!\left(h_k,G_e(m_k)u_k^{ext},G_i(m_k)u_k^{int},G_z(m_k)z_k,\Theta_k;\Delta t_k\right),
\qquad
\tilde a_k=G_a(m_k)a_k .
\]
各增益是有单位接口上的线性或非线性算子，不能把它们相加为一个抽象温度。模型还需声明采样步长、通道延迟、噪声和模式切换条件；若模式在一次步长内变化，应使用事件重置而非假定连续可导。

入睡后外部校正通常减弱，内部模型相对主导，于是生成内容更容易被当前系统当作场景证据。但“相对主导”不等于因果律失效。梦中对象仍由既有记忆、生成模型、情绪和身体信号约束；门铃、疼痛或呼吸困难也可能被编入梦境并触发醒来。系统应区分真实外部消息、内部重放、模型补全和梦中再解释，若醒后报告无法恢复来源，就只能将内容标记为低来源置信的体验报告，而不能据此重写外部事实。

感知边界的重配置可以用校准比描述。给定任务读出 $R_q$，定义外部证据敏感度
\[
\chi_e=\left\|\frac{\partial R_q(h_{k+1})}{\partial u_k^{ext}}\right\|,
\qquad
\chi_z=\left\|\frac{\partial R_q(h_{k+1})}{\partial z_k}\right\|.
\]
当 $\chi_z/\left(\chi_e+\epsilon\right)$ 增大时，报告更受内部生成影响；该比值依赖查询、通道和状态，不是梦境强度的普适标量。若外部刺激足够强，$\chi_e$ 仍可能上升并改变模式。验证需要在不同睡眠阶段施加受控声音、气味或触觉，比较报告、神经指标和觉醒概率，而不是从方程名称直接推出现象。

动作边界同样是门控而非绝对切断。快速眼动睡眠中的肌张力抑制能降低多数随意动作输出，但眼动、呼吸和局部肌肉活动仍存在；不同睡眠状态和个体也有差异。工程系统可通过影子执行器实现类似隔离：计划器产生动作候选，模拟环境返回虚拟回执，真实执行器保持只读或需要额外授权。此时必须给虚拟回执标记来源，防止离线模拟被写成真实经历。若动作门控异常开启，系统还要有硬件边界和紧急停止，不能只依赖上层意图。

边界变化会改变可辨识性。外部反馈减少时，多种内部模型可能生成相似梦境，单次报告无法判断是哪一机制主导；醒后叙述又受到遗忘和语言重构影响。我们因此把梦中状态、觉醒瞬间记录、延迟报告和后来解释分层保存。对人类研究，主观报告、行为、生理信号和神经数据各提供有限证据；对AgentOS离线维护，则可直接记录生成轨迹、模型版本、外部输入与执行门控，但即便日志完整，也不能由功能相似推出主体体验相同。

本节保留“从处理现实转向模拟现实”的直觉，但将其限定为输入—生成权重与行动权限的重新配置。验收使用残余外部刺激、内感受扰动、模拟回执污染、动作门控故障和模式快速切换。系统应识别消息来源，避免把模拟写成事实，在外部风险上升时能够唤醒或升级控制，并在恢复清醒闭环后重新校准环境版本。报告外部敏感度、生成敏感度、来源误判、动作泄漏、唤醒延迟和实际任务结果；这些指标支持离线生成模型，不证明睡眠是绝热演化或思维波在边界全反射。
